\chapter{Amplitude, Energy, Saturation, and Pedestal}

\section{Scope and raw-data gate}

This chapter collects the studies concerned with pulse height, integrated
charge, saturation recovery, energy ordering, and baseline estimation.  The
common starting point is the S00 B-stack selected-pulse population rebuilt from
raw ROOT files, using the median of samples 0--3 as the local baseline and a
baseline-subtracted even-channel peak threshold of 1000 ADC on B2/B4/B6/B8.
The studies in this chapter repeatedly reproduce the same entry count,
\[
  N_{\mathrm{sel}}=640737,
\]
before fitting any calibration, regression, or closure model.  That raw
reproduction gate matters because most targets in this chapter are proxy
targets: an odd duplicate readout, an artificially clipped clean pulse, a
downstream charge proxy, or a leave-one-pre-trigger-sample closure target.
They are useful tests of method performance, but they are not direct
per-event deposited-energy truth.

The primary resolution metric for amplitude, charge, and energy-proxy rows is
the 68th percentile absolute fractional residual,
\[
  r_{68}
  = Q_{0.68}\left(\left|\frac{\hat{y}_i-y_i}{y_i}\right|\right),
\]
evaluated on held-out runs or leave-one-run-out folds.  Pedestal studies use
the held-out mean absolute error,
\[
  \mathrm{MAE}
  = \frac{1}{N}\sum_i |\hat{p}_i-p_i|,
\]
where \(p_i\) is an excluded pre-trigger sample.  The quoted intervals are
bootstrap confidence intervals from the corresponding held-out records or
run-blocks.  All method comparisons below are therefore conditional on the
available beam-triggered mirror and on the chosen proxy target.

\section{Duplicate-readout amplitude and charge closure}

The P04 family tests whether the even-channel waveform can predict an
independent odd-channel duplicate readout.  For an event \(i\), the closure
target is
\[
  A_i^{\mathrm{odd}} = \max_t[-w^{\mathrm{odd}}_{it}],
  \qquad
  Q_i^{\mathrm{odd}} = \sum_t \max[-w^{\mathrm{odd}}_{it},0],
\]
while the model inputs are restricted to the even-channel waveform and derived
even-channel summaries.  Run identifiers, event identifiers, and odd-channel
samples are excluded.  This makes the task a strong electronic/readout closure
benchmark, not an absolute energy calibration.

The traditional baselines are per-stave log-linear calibrations of even peak
or even positive-lobe integral to the odd duplicate readout, plus shifted
template-scale fits.  The strongest traditional repair in P04d augments direct
template diagnostics with peak, charge, and shape summaries in a Huber
calibrator.  The ML methods are histogram-gradient-boosted trees or
ExtraTrees regressors trained on the 18 waveform samples and non-forbidden
shape summaries.  Held-out runs 57 and 65 are absent from training.

\begin{table}[H]
\centering
\caption{Duplicate-readout closure on held-out runs.  Lower \(r_{68}\) is
better.  The P04 rows use odd duplicate amplitude or charge as target; P04d is
the strengthened amplitude closure benchmark.}
\begin{tabular}{llccc}
\toprule
Study & Method & Target & \(r_{68}\) & 95\% CI \\
\midrule
P04 & peak calibrated & amplitude & 0.1238 & [0.1221, 0.1251] \\
P04 & HGB regressor & amplitude & 0.0091 & [0.0090, 0.0093] \\
P04 & integral calibrated & charge & 0.1954 & [0.1936, 0.1975] \\
P04 & HGB regressor & charge & 0.0151 & [0.0148, 0.0153] \\
P04d & strong Huber traditional & amplitude & 0.0203 & [0.0199, 0.0206] \\
P04d & ExtraTrees regressor & amplitude & 0.0027 & [0.0026, 0.0028] \\
\bottomrule
\end{tabular}
\label{tab:amplitude-closure}
\end{table}

The winner for duplicate-readout closure is the ML waveform model.  In the
original P04 benchmark the HGB amplitude row improves \(r_{68}\) by
0.1147 absolute fractional units relative to the peak baseline, and the charge
row improves by 0.1803.  P04d shows that a much stronger traditional Huber
calibrator closes much of the gap, but the ExtraTrees model remains lower by
0.0176 in \(r_{68}\).  The direct shifted-template scale itself performs
poorly: changing the fit window, adding a baseline nuisance, peak-normalising
the template, or switching to Huber loss does not repair the direct-scale
pathology.  The useful conventional repair is post-fit calibration with
diagnostics, not the raw template scale.

The leakage audits are also part of the result.  The duplicate target rows
remove invalid odd targets after reproduction, enforce train/held-out run
separation, check zero train/test event-key overlap, exclude forbidden odd
samples, and compare to shuffled-target and stave-only sentinels.  The very
small ML errors are therefore plausible as duplicate-readout closure, but they
must not be promoted to a deposited-energy resolution without an external
energy truth.

\section{Saturation recovery}

Saturation recovery asks whether an amplitude lost to clipping can be
recovered from the surviving rising edge.  P07 constructs a leakage-free
self-truth benchmark by taking clean unsaturated pulses with known amplitude,
applying a constant artificial ADC ceiling \(C\), and asking methods to
recover the original amplitude.  The constant ceiling is essential.  An earlier
ceiling proportional to the true amplitude leaked the target through the
clipped maximum; the corrected benchmark fixes \(C\), so the saturated maximum
is no longer an amplitude label.

The traditional method fits a mean clean-pulse template to the unclipped
rising-edge samples.  The ML method is a gradient-boosted regressor on the
clipped waveform.  The split is by run: runs 58--61 for training and
62, 63, and 65 for testing.

\begin{table}[H]
\centering
\caption{P07 artificial-saturation amplitude recovery on held-out runs.
Entries are \(r_{68}(|\hat{A}-A|/A)\).}
\begin{tabular}{ccccc}
\toprule
Ceiling \(C\) [ADC] & Test pulses & Naive ceiling & Template & GBR \\
\midrule
4000 & 8873 & 0.264 & 0.104 & 0.032 \\
3000 & 20254 & 0.346 & 0.239 & 0.039 \\
2500 & 27971 & 0.403 & 0.233 & 0.042 \\
2000 & 33823 & 0.493 & 0.286 & 0.046 \\
\bottomrule
\end{tabular}
\label{tab:saturation-recovery}
\end{table}

\begin{figure}[H]
\centering
\includegraphics[width=0.82\textwidth]{\detokenize{../../reports/P07_saturation_recovery/fig_saturation_recovery.png}}
\caption{P07 saturation recovery benchmark.  The ML regressor degrades
slowly as the artificial ceiling is lowered, whereas the template extrapolation
degrades rapidly.}
\label{fig:saturation-recovery}
\end{figure}

The winner for the artificial saturation benchmark is ML: the GBR recovers
amplitude at about 3--5\% \(r_{68}\), compared with 10--29\% for the template
extrapolation.  The likely reason is that the rising edge is not a fixed
linear scale copy of one template; it carries amplitude-dependent shape
information that tree regressors exploit.  The systematics are nevertheless
substantial.  The benchmark uses an ideal hard clip, real saturated pulses have
no truth label, and natural high-B2 pulses may include late components or
overlaps absent from the clean-pulse training population.  Later P07 boundary
studies therefore keep the result as a strong recovery method under artificial
truth, not as an unconditional production correction.

\section{Energy proxy and range-ordering limits}

The S14 studies propagate the charge and saturation results into range-energy
preflight tests.  The geometry/PSTAR information is used only as a monotonic
ordering envelope.  It supplies a physically motivated ordering check, not
particle-identification truth and not an absolute per-event deposited energy.

In S14c, the nominal 4 cm geometry comparison uses observed even charge,
a traditional rising-edge saturation correction, and an ML P07/P04 corrected
charge model.  The ML model excludes run, event, depth, PID, and odd samples.
It reduces the internal energy-proxy \(r_{68}\) from 0.0289 for the traditional
template-corrected row to 0.0145, with no depth-order violations in the tested
envelope.  However, the external charge-propagated follow-up S14c topology
study shows that downstream charge proxies do not reach the 10\% target:
the best combined row is about 0.178 with a 95\% CI of [0.169, 0.198].

\begin{table}[H]
\centering
\caption{Energy-related proxy rows.  The first block is the S14c internal
ordering proxy at nominal 4 cm geometry; the second block is the downstream
charge-propagated preflight.}
\begin{tabular}{llccc}
\toprule
Study & Method & Quantity & \(r_{68}\) & 95\% CI \\
\midrule
S14c & observed even charge & energy proxy & 0.0212 & [0.0198, 0.0225] \\
S14c & traditional template corrected & energy proxy & 0.0289 & [0.0248, 0.0350] \\
S14c & ML P07/P04 corrected & energy proxy & 0.0145 & [0.0142, 0.0150] \\
S14c topology & best ML charge-propagated row & combined proxy & 0.1781 & [0.1692, 0.1977] \\
S14c topology & best traditional propagated row & combined proxy & 0.2015 & [0.1875, 0.2156] \\
\bottomrule
\end{tabular}
\label{tab:energy-proxy}
\end{table}

The winner for internal charge/ordering diagnostics is again ML, but the
physics conclusion is a limitation: the data-only analysis does not provide
absolute energy resolution at the 10\% level.  The strong duplicate-readout
closure and artificial saturation recovery do not survive unchanged when the
target becomes downstream topology-conditioned charge.  The remaining spread
is dominated by support and topology variability rather than by an obvious
hidden B2 waveform regressor failure.

\section{Pedestal and baseline validation}

Pedestal studies test whether the adaptive positivity-constrained baseline is
unbiased.  Because the current mirror contains no true forced/random pedestal
sample, S16 and S16b use leave-one-pre-trigger-sample closure.  For each pulse
and each early sample \(k\in\{0,1,2,3\}\), the estimator must predict
\(w_k\) without seeing that sample.  The adaptive positivity correction lowers
the leave-one-out median only enough to satisfy
\[
  \min_{t \in \mathcal{S}_{\mathrm{nonjagged}},\,t\ne k}
  \left(w_t-\hat{p}_{-k}\right) \geq -\epsilon(A),
  \qquad
  \epsilon(A)=\max(25\ \mathrm{ADC},0.015A).
\]
That construction can guarantee non-negative corrected samples, but the
guarantee is not an independent validation of the pedestal.

S16 compares median/mean pre-trigger estimators, the adaptive positivity
correction, and a calibrated histogram-gradient-boosted regressor.  S16b then
strengthens the traditional set with a line-through-three early-sample
predictor and calibrated low-sample estimators; its ML method is a regularized
ridge closure model.

\begin{table}[H]
\centering
\caption{Pedestal closure on held-out runs.  Lower MAE is better.  S16b is the
leakage-aware follow-up with stronger traditional estimators.}
\begin{tabular}{llcc}
\toprule
Study & Method & MAE [ADC] & 95\% CI [ADC] \\
\midrule
S16 & mean3 & 260.70 & [236.25, 287.99] \\
S16 & median3 & 273.64 & [244.24, 302.67] \\
S16 & adaptive positivity correction & 341.04 & [300.45, 373.27] \\
S16 & calibrated HGBR & 48.88 & [43.82, 55.29] \\
S16b & line3 predictor & 169.34 & [135.28, 200.24] \\
S16b & calibrated ridge closure & 173.71 & [149.60, 198.53] \\
S16b & calibrated low2 traditional & 211.42 & [167.23, 264.49] \\
\bottomrule
\end{tabular}
\label{tab:pedestal-closure}
\end{table}

\begin{figure}[H]
\centering
\includegraphics[width=0.82\textwidth]{\detokenize{../../reports/1780997954.15337.77205a71__s16_pedestal_baseline_validation/fig_heldout_residual_distributions.png}}
\caption{S16 held-out pre-trigger residual distributions for simple,
adaptive, and ML pedestal estimators.}
\label{fig:pedestal-residuals}
\end{figure}

\begin{figure}[H]
\centering
\includegraphics[width=0.82\textwidth]{\detokenize{../../reports/1781000826.539659.030b7796__s16b_independent_pedestal_estimator_closure/fig_head_to_head.png}}
\caption{S16b leakage-aware head-to-head pedestal closure.  The strengthened
line3 traditional estimator and ridge ML closure are statistically tied in
MAE on the held-out records.}
\label{fig:pedestal-head-to-head}
\end{figure}

The raw S16 ML row is numerically the best leave-one-sample predictor, but the
accepted pedestal conclusion is more conservative.  In the leakage-aware S16b
benchmark, the best traditional method is the line3 predictor at
169.34 ADC MAE, while calibrated ridge gives 173.71 ADC with a paired
ML-minus-traditional delta of 4.38 ADC and CI [-15.61, 19.58].  Thus the
accepted winner for an operational pedestal estimator is traditional or tied,
not ML.  ML remains useful as a contamination diagnostic because later waveform
samples can predict an early contaminated sample, but that predictive closure
does not prove an unbiased zero-signal electronics pedestal.

\section{Systematics and caveats}

\begin{table}[H]
\centering
\caption{Dominant systematic limitations by topic.}
\begin{tabular}{p{0.20\textwidth}p{0.35\textwidth}p{0.35\textwidth}}
\toprule
Topic & Main systematic & Consequence \\
\midrule
Duplicate amplitude/charge & Odd duplicate readout is an electronic closure target, not deposited-energy truth. & ML wins closure, but the result cannot by itself calibrate absolute energy. \\
Template scale & Direct shifted-template scale is pathological under multiple windows and losses. & Strong traditional comparison must use calibrated diagnostics, not raw template scale. \\
Saturation & Artificial hard clipping gives self-truth; real saturated pulses have no amplitude truth and may have overlaps. & ML is preferred for artificial recovery, while natural transfer needs boundary and timing-tail audits. \\
Energy proxy & PSTAR/depth information supplies only monotonic ordering; downstream topology support is sparse. & Internal ordering improves with ML, but the 10\% absolute-energy preflight fails. \\
Pedestal & No true forced/random pedestal sample is present in the current mirror. & Pre-trigger closure can diagnose contamination, but it is not a no-pulse pedestal measurement. \\
\bottomrule
\end{tabular}
\label{tab:amplitude-systematics}
\end{table}

\section{Chapter verdict}

The robust pattern is that waveform ML wins when the target is an independent
shape-rich closure target: duplicate-readout amplitude/charge closure and
artificial saturation recovery.  The strongest numbers are P04d ExtraTrees
amplitude closure at \(r_{68}=0.0027\) versus 0.0203 for a calibrated Huber
traditional method, and P07 GBR saturation recovery at 0.032--0.046 versus
0.104--0.286 for template extrapolation.  The energy chapter is nevertheless
bounded by truth availability.  The S14 proxy rows show improved internal
ordering with ML, but downstream charge-propagated energy remains at about
18--20\% and fails the 10\% target.  Pedestal validation is even more
truth-limited: without forced/random triggers, the defensible operational
conclusion is that strong traditional early-sample closure ties ML, while ML
serves as a sensitive diagnostic for pre-trigger contamination.
